\section{The discretisation}
\label{sec:disc}

The operator is a control-volume finite-element (CVFEM) discretisation of incompressible
Navier--Stokes on trilinear hexahedra. Each of the eight nodes of an element owns one
sub-control volume, and the element contributes to the flux balance of those volumes across
twelve \emph{sub-control surfaces} (SCS). The unknowns are collocated: each node carries
$(u_x,u_y,u_z,p)$, so an element holds $8\times 4 = 32$ degrees of freedom.

\subsection{Sub-control surfaces}

The twelve surfaces are enumerated in \code{CVFEM\_HEX8\_SCS}
(\code{src/upwind/cvfem\_hex8\_ns\_upwind\_kernels.hpp}). Each entry is a triple
$(i, j, \Avec^{\mathrm{ref}})$: the two nodes whose sub-control volumes the surface separates,
and the surface's area vector on the reference element. The enumeration is grouped
\emph{four per reference direction}:

\begin{center}
\begin{tabular}{@{}llll@{}}
\toprule
direction & $\scs$ & node pairs $(i,j)$ & $\Avec^{\mathrm{ref}}$ \\
\midrule
$\xi$   & $0,1,2,3$    & $(0,1)\ (3,2)\ (4,5)\ (7,6)$ & $(\tfrac14,0,0)$ \\
$\eta$  & $4,5,6,7$    & $(0,3)\ (1,2)\ (4,7)\ (5,6)$ & $(0,\tfrac14,0)$ \\
$\zeta$ & $8,9,10,11$  & $(0,4)\ (1,5)\ (2,6)\ (3,7)$ & $(0,0,\tfrac14)$ \\
\bottomrule
\end{tabular}
\end{center}

Two facts follow from this table and are used throughout the implementation. First,
$\scs \,/\, 4$ is the reference direction, and because $\scs$ is a template parameter in every
face kernel the direction is a compile-time constant. Second, each pair $(i,j)$ is
\emph{edge-adjacent}: the two nodes differ in exactly one logical coordinate, the one naming the
group. Section~\ref{sec:geom} shows that this is what makes the element Jacobian sufficient to
describe the surface geometry on the affine path.

Every node appears in exactly three of the twelve surfaces, once per direction group. That
multiplicity is the reason the element's four accumulators per node ($r_x, r_y, r_z$ and the
continuity row) are live across the whole surface loop, which
Section~\ref{sec:mk:structure} returns to as the binding constraint on the kernel.

\subsection{The residual}

Writing $R_a$ for the contribution to node $a$'s sub-control volume, the element residual is a
sum over the twelve surfaces of a flux times an area, plus the volume terms:
\begin{equation}
R_a \;=\; \sum_{\scs \,:\, i(\scs)=a} \mathbf{F}_\scs
      \;-\; \sum_{\scs \,:\, j(\scs)=a} \mathbf{F}_\scs
      \;+\; \text{(transient)}_a \;+\; \text{(boundary)}_a .
\label{eq:residual}
\end{equation}
The antisymmetry is exact and structural: a surface adds its flux to node $i$ and subtracts the
same value from node $j$, so the element conserves by construction regardless of what
$\mathbf{F}_\scs$ is. Every kernel in this document preserves it literally --- the flux is
computed once per surface and applied with opposite signs --- which is why no scheme described
here can break conservation, only accuracy.

The surface flux has three parts, treated in Section~\ref{sec:components}: a viscous part, a
convective part carrying the mass flux, and the pressure--velocity coupling that fixes the mass
flux itself. The convective part is where the schemes differ, and Section~\ref{sec:schemes}
takes those one at a time.
